\documentclass[12pt]{article}
\usepackage[utf8]{inputenc}
\usepackage[T1]{fontenc}
\usepackage{hyperref}
\title{Buhi'non Bikol Language Resources for K-12 Curriculum Development}
\author{Compiled by Claude Sonnet 5 (claude-sonnet-5) for Andr\'as Kornai}
\date{}
\begin{document}
\maketitle

\section{Introduction}
Buhi'non Bikol (Buhinon) is an Austronesian language of the Bikol group. It is spoken in Buhi, Camarines Sur, Philippines. Wikipedia says it has features of both the Polangui (Albay) Bikol and the Iriga variant of Rinconada Bikol. Global Recordings Network lists it as ISO 639-3 ubl, with varieties such as Tambo' and San Vicente. One database (ABVD) notes a different code, bhk, so coding is inconsistent across sources. Very little educational material turned up in the searches.

\section{Primary Educational Resources}
No primers, textbooks or school materials in Buhi'non turned up in the searches.

\section{Linguistic and Reference Materials}
\begin{itemize}
\item Wikipedia, Bikol languages (classification and a brief description of Buhinon): \url{https://en.wikipedia.org/wiki/Bikol_languages}
\item Austronesian Basic Vocabulary Database, Bikol (Buhinon) wordlist, citing a 2009 wordlist by Olson, Ballenas and Borromeo recorded in 2006: \url{https://abvd.eva.mpg.de/austronesian/language.php?id=710}
\end{itemize}

\section{Religious and Community Resources}
\begin{itemize}
\item Joshua Project profile, Buhi'non Bikol in the Philippines (mission-oriented): \url{https://joshuaproject.net/people_groups/20528}
\end{itemize}

\section{Digital and Audio Resources}
\begin{itemize}
\item Global Recordings Network, Bikol, Buhi'non (audio recordings page for ubl; contents not inspected): \url{https://www.globalrecordings.net/en/language/27531}
\end{itemize}

\section{Curriculum Development Context}
The reference description says the language is not known to be taught in schools. The searches found no orthography, grammar or textbook. In the Philippines, a curriculum would likely build on the national mother-tongue-based multilingual education framework, but the searches found no document on that framework. Related Bikol varieties with larger bodies of material could be a starting point. Their suitability for Buhi'non is not established by these searches.
\end{document}